\documentclass[10pt]{article}
\usepackage[margin=20mm]{geometry}
\usepackage{amsmath,amssymb,fontspec,unicode-math,xurl,hyperref}
\IfFontExistsTF{DejaVu Serif}{\setmainfont{DejaVu Serif}}{\setmainfont{Latin Modern Roman}}
\setmathfont{latinmodern-math.otf}
\setlength{\parindent}{0pt}\setlength{\parskip}{6pt}\sloppy
\begin{document}
Abstract: Three defects affect the stated theoretical guarantees: a purported linear spline construction is quadratic, a clipping threshold depending on the current response invalidates a conditional expectation identity, and the stated rate vanishes on a permitted zero-variation signal. Explicit counterexamples establish these claims. They support Major assessments of the affected derivations and theorem statements, without establishing that the intended positive-variation rate or the empirical advantages are false.\par
Author: AI-assisted technical review\par
Date:\par
Title: Technical Verification of AKORN\par
\# Identity and scope\par
The reviewed paper is Sunil Madhow, Dheeraj Baby and Yu-Xiang Wang, AKORN: Adaptive Knots generated Online for RegressioN splines, ICML 2025, PMLR 267, OpenReview `vfRmp9RPV7` (0292/I25-43661) [Retained latex: \textbackslash{}cite\{paper\}]. The source is the frozen proceedings PDF; one-based PDF pages coincide with printed page numbers in the cited locations. Released code was inspected separately at `SunilMadhow/AKORN`, commit `13bb80f30c1e212b4c6c2a0d7453aba87a70bcc6` [Retained latex: \textbackslash{}cite\{code\}]. Its conference-time provenance is not established.\par
The checks cover the principal algorithm definitions, Theorems 6.1–6.2, the projection argument in Lemma C.2, the knot-count/projection proof chain in Appendix C, and the conditional-noise step in Appendix D. They include selected static code correspondence and exact independently written calculations. Appendix E’s full approximation analysis, Appendix F’s nonuniform-design extension, external cited proofs and empirical reproduction were not exhaustively verified. No author code was executed.\par
\# Definitions and the core dependencies\par
The observations are \(Y_i=f(x_i)+\epsilon_i\), with \(x_i=i/n\), independent \(\epsilon_i\sim N(0,\sigma^2)\), bounded \(|f|\le1\), and \[C=\sum_{i=3}^n\left|
 \frac{f(x_i)-f(x_{i-1})}{x_i-x_{i-1}}-
 \frac{f(x_{i-1})-f(x_{i-2})}{x_{i-1}-x_{i-2}}\right|<\infty .\] The paper permits \(C=0\). Let \(S(K)\) be continuous linear splines evaluated on the design, with knots \(K\), and \(F(K)\) the larger piecewise-linear space permitting discontinuities. AKORN outputs the least-squares projection \(P_{S(K)}Y\). Lemma C.7 supplies a spline sandwiched between forward and backward fits to control approximation error. Lemma D.6 supplies the noise identity used in the martingale argument for ADDLE; ADDLE’s guarantee in turn supports knot detection. Both dependencies are central to Theorem 6.2.\par
\# Substantive findings\par
Equation reference: eq:bad. Problem 1 (The constructed function in Lemma C.7 is not a linear spline).  Lemma C.7, Case 2 (p. 18), sets \[
 s(z)=g(z)+(h(z)-g(z))\frac{z-a}{b-a},
 \qquad a=k_{i-1},\ b=k_i,\] and asserts membership in the linear-spline space. The claimed implication from noncrossing to linearity is false.\par
\#\# Complete check\par
Use \(n=8\), \(K_f=\{1/2\}\), \(K_b=\{1/4\}\), and, on \([0,1]\), \[g(z)=\max\{0,z-1/2\},\qquad h(z)=\max\{5/4,1+z\}.\] These are in \(F(K_f)\) and \(F(K_b)\), respectively; continuous piecewise-linear functions are allowed members. The knot sets are disjoint, and \(g\le1/2<5/4\le h\), so there are no crossover points, regardless of whether crossings are tested on the grid or continuously. The first knot \(1/4\) belongs to \(K_b\), so the proof’s base case gives \(s=g=0\) up to that knot. The next knot \(1/2\) belongs to \(K_f\), satisfying precisely the hypotheses of Case 2.\par
On \([1/4,1/2]\), Equation [eq:bad] (\#eq:bad) becomes \[s(z)=4z^2+3z-1,\qquad s''(z)=8.\] All three of \(1/4,3/8,1/2\) are design points, but \[s(1/4)=0,\qquad s(3/8)=11/16,\qquad s(1/2)=3/2.\] An affine segment with these endpoints must have midpoint \(3/4\). The discrepancy is \(-1/16\). Thus membership fails even if \(S(K)\) is interpreted only as vectors of evaluations on the finite design. Exact rational arithmetic reproduces these values in `verification/0292/check.json`. The original p. 18 was rendered and inspected to exclude extraction error.\par
\#\# Counterevidence, correction and impact\par
The pointwise convex combination in Equation [eq:bad] (\#eq:bad) does stay between \(g\) and \(h\); that fact does not make it affine. Equal slopes would be a valid special case, but equal slopes are not a lemma hypothesis. The lemma’s existence conclusion is not disproved by this example. Indeed, the local endpoint interpolant \[s_*(z)=6z-\tfrac32\quad (1/4\le z\le1/2)\] works here: \(s_*\ge g=0\) and \(h-s_*=5/2-5z\ge0\). More generally, endpoint interpolation stays in an ordered strip between two affine functions because the relevant differences are affine and nonnegative at both endpoints. This verifies a local repair, not all jump, crossover and collision cases in the full construction.\par
Theorem C.9 (pp. 18–19), hence Theorem 6.2, invokes Lemma C.7 to produce a member of \(S(K)\) with controlled error. The written construction fails at that exact required membership step. This is a Major defect in a core derivation even if a corrected construction can restore the lemma; it is not a proof that the desired rate is unattainable.\par
Problem 2 (The clipping threshold destroys the conditional-noise identity).  Algorithm 5 (p. 13) uses \[B=\max_{1\le i\le n}|Y_i|+
 \max\{\sigma\sqrt{2\log(4n/\delta)},1\}\] to clip the linear experts. Lemma D.6(1) (p. 20), with \(\mathcal F_j=\sigma(Y_1,\ldots,Y_{j-1})\), asserts \[E[(Y_j-z_j)^2-(Y_j-\theta_j)^2\mid\mathcal F_j,V]
 =E[(z_j-\theta_j)^2\mid\mathcal F_j,V].\] The identity fails because \(B\), and consequently \(z_j\), depends on the current response.\par
\#\# Parameters, hypotheses and calculation\par
Take \(n=3\), \(\sigma=0.1\), \(\delta=0.1\), and \((\theta_1,\theta_2,\theta_3)=(-1,1,1)\). A permissible underlying function equals \(-1\) up to \(x_1\), rises linearly with slope \(6\) to \(x_2\), and then equals \(1\). Thus it is bounded by one, has the required design values and \(C=6\).\par
Equation reference: eq:noise. Write \(a=\sigma\sqrt{2\log120}=0.3094347<1\) and \(V=\{|\epsilon_i|\le a\text{ for all }i\}\). First evaluate the conditional law at past responses \(Y_1=-1,Y_2=1\); their noise values satisfy \(V\). The linear expert starting at the first observation extrapolates to \(3\) at \(x_3\). Conditional on the past and \(V\), \(e=\epsilon_3\) has a symmetric normal law truncated to \([-a,a]\), and \[B=2+\max(e,0)<3,\qquad z_3=B.\] All clipping operations are therefore in the stated clipped branch. Subtracting the right side of Lemma D.6(1) from its left side gives \[\begin{aligned}
 -2E[e(z_3-\theta_3)\mid V,Y_1=-1,Y_2=1]
 &=-2E[e(1+\max(e,0))\mid |e|\le a]\notag\\
 &=-E[e^2\mid |e|\le a]<0. \end{aligned}\] The exact conditional moment has the analytic form \[v=\sigma^2\left(1-\frac{2c\phi(c)}{2\Phi(c)-1}\right),
 \qquad c=a/\sigma,\] and independent evaluation gives \(v=0.0097938491\). Thus the claimed zero difference is approximately \(-0.0097938491\).\par
The exact past point has probability zero, so it is not itself asserted to be a positive-probability counterevent. Conditional moments and this clipped prediction vary continuously with the two past responses in a small neighborhood of \((-1,1)\), while \(V\) has a positive-density interior there. The strictly negative difference in Equation [eq:noise] (\#eq:noise) therefore persists on a neighborhood of positive probability. This establishes failure of the asserted conditional identity as an almost-sure statement.\par
\#\# Alternatives and dependence\par
Conditioning on \(V\) alone preserves the symmetric truncated law of each independent noise variable. The failure comes from the response-dependent threshold, not from an automatic claim that every conditioning argument is invalid. The counterexample follows Equation (1)’s Gaussian model and Section 3.1’s \(|f|\le1\) bound. Algorithm 3’s \([0,1]\) history annotation conflicts with that model; imposing it would exclude theorem-covered data. A deterministic bound or a threshold measurable with respect to the past would remove this particular dependence, but would change the algorithm or its assumptions and require a fresh proof. The released code’s omission of this clipping step is additional evidence that its guarantees cannot simply be equated with the printed algorithm.\par
Lemma D.8 (p. 21) uses Lemma D.6 to declare the corresponding noise term a martingale difference and apply a Freedman-type inequality. The counterexample invalidates that justification. This is Major for the ADDLE error proof and the dependent AKORN guarantee. No failure rate of the empirical estimator is inferred; no complete replacement concentration argument has been verified.\par
Problem 3 (The stated bounds incorrectly vanish on affine signals).  Theorems 6.1 and 6.2 (p. 7) quantify over \(C<\infty\), but give total error \(\widetilde O(n^{1/5}C^{2/5})\) and mean squared error \(\widetilde O(n^{-4/5}C^{2/5})\), respectively, without a parametric baseline.\par
For Theorem 6.2 choose \(f=0\), any \(n\ge3\) and \(\sigma>0\). All stated assumptions hold and \(C=0\). For every selected knot set, \(S(K)\) contains constants. Least-squares orthogonality consequently gives, pathwise, \(\overline{\widehat f}=\overline Y=\overline\epsilon\), even though \(K\) is data-dependent. Therefore \[\frac1n\|\widehat f\|_2^2\ge\overline\epsilon^{\,2}>0
 \quad\text{almost surely},\qquad
 \frac{n\overline\epsilon^{\,2}}{\sigma^2}\sim\chi^2_1.\] The displayed upper bound equals zero; it cannot hold with probability \(1-\delta>0\). For the online claim, the constant signal \(f=1/2\) also has \(C=0\), whereas the empty-history first prediction is zero, already giving total error at least \(1/4\).\par
Multiplicative constants and logarithms cannot turn zero into a positive upper bound. There is no stated lower bound on \(C\). Counterevidence for the intended interpretation is explicit: Lemma C.1 retains \(\max\{n^{1/5}C^{2/5},1\}\), and Section 5 discusses an \(n^{-1}\) parametric rate. Thus a parametric term was plausibly intended. Nevertheless, it disappears in the theorem statements and later Corollary C.6/C.9 bounds. Retaining a baseline such as \(\widetilde O(n^{-1}+n^{-4/5}C^{2/5})\) is necessary for a uniform offline statement. This report does not certify that amended bound, because Problems 1 and 2 leave other proof obligations open. The assessment is Major for the stated central uniform guarantee; the positive-\(C\) leading exponent is not refuted by this check.\par
\# Source correspondence and surviving results\par
Lemma C.2’s nested-projection argument is sound: since \(S\subset F\), \(\|P_SY-P_S\theta\|\le\|P_FY-\theta\|\), and the triangle inequality gives the displayed factor-two bound after squaring. This deterministic statement remains valid for a selected random \(K\). The interval-uniform strategy in Lemma C.3 also avoids simply treating adaptive knots as independent of the responses; that idea is not the source of Problem 2.\par
The released `akorn.py`, lines 71–81, uses continuous `RegressionSpline` preliminary fits rather than the \(F(K)\) fits in Algorithm 2; `splines.py`, lines 25–29, confirms the basis \(1,x,(x-k)_+\). It also omits collision-repair knots, a proof requirement acknowledged in Remark 4.2, and uses a pseudoinverse. The historical-indexed predictions in `flh.py` agree with Appendix C’s intended stopping statistic rather than a literal repeated latest-prediction subscript in Algorithm 1. These counterexamples are therefore not attributed indiscriminately to every release branch. The exact theoretical coverage of this different released estimator and its connection to the reported runs remain Needs Evidence.\par
The reported simulations, visual fit quality and possible existence of an optimal repaired estimator can survive the demonstrated proof defects. No empirical scores were rerun or revised. Later versions and author errata were not supplied or separately adjudicated.\par
\# Conclusion\par
Item | Recorded state\par
C.7 construction | Major: written function leaves the required spline space\par
D.6 noise identity | Major: explicit nonzero conditional cross term\par
Zero-variation theorem domain | Major: stated zero error contradicts permissible signals\par
Repair verification | Local interpolation checked; no complete repaired theorem\par
Released estimator’s rate | Needs Evidence\par
Scientific adjudication | Pending\par
Author adoption / human review | Unknown, not requested / none recorded\par
9 S. Madhow, D. Baby and Y.-X. Wang. AKORN: Adaptive Knots generated Online for RegressioN splines. PMLR 267, ICML 2025. Frozen proceedings version. https://proceedings.mlr.press/v267/madhow25a.html. `SunilMadhow/AKORN`. Commit `13bb80f30c1e212b4c6c2a0d7453aba87a70bcc6`. https://github.com/SunilMadhow/AKORN.\par
\# Reproducibility and limitations\par
The original PDF SHA-256 is . All staged input hashes and preserved archive correspondence are in `inputs.json` and the unchanged source manifest. Within `verification/0292`, `check.py/check.json` contain the exact rational construction, and `check2.py/check2.json` contain the independent truncated-normal calculation. These are algebraic and finite arithmetic diagnostics, not execution of author software. Auxiliary notes `a1.json` and `a2.json` retain source and proof checks. The primary assistant independently checked the derivations, original formula images, exact TeX build and report pages; native auxiliaries performed source inspection and an independent proof challenge. No human review or author adoption is represented. Build, visual and source evidence are hash-bound in the corresponding JSON records.\par
\end{document}
